% © 2026 Ing. David Jaroš — CC BY-NC-ND 4.0
%
% This work is licensed under a Creative Commons Attribution-NonCommercial-NoDerivatives
% 4.0 International License (CC BY-NC-ND 4.0).
%
% License History: Earlier drafts (up to v0.3) were released under CC BY 4.0.
% From v0.4 onward, all material is released under CC BY-NC-ND 4.0 to protect
% the integrity of the theoretical work during ongoing academic development.

% UBT-AI-PROVENANCE-BEGIN
% schema: ubt-ai-provenance/v1
% tier: B_machine_verified
% ai_assistance: disclosed
% human_review: machine-verification
% editorial_responsibility: Ing. David Jaroš
% policy: ../../AI_PROVENANCE.md
% notice: Machine-verified against named sources or verifiers; individual attestation is not claimed.
% UBT-AI-PROVENANCE-END

\documentclass[12pt,a4paper]{article}
\usepackage{amsmath,amssymb,amsthm}
\usepackage{mathtools}
\usepackage{geometry}
\usepackage{booktabs}
\usepackage{hyperref}
\usepackage{xcolor}
\usepackage{tcolorbox}
\tcbuselibrary{skins,breakable}
\geometry{left=2.5cm,right=2.5cm,top=2.5cm,bottom=2.5cm}

\newtheorem{theorem}{Theorem}[section]
\newtheorem{lemma}[theorem]{Lemma}
\newtheorem{remark}[theorem]{Remark}
\newtcolorbox{statusbox}[1][]{colback=gray!8!white,colframe=black!60,arc=3pt,boxrule=1pt,title=Status Summary,#1}

\newcommand{\Neff}{N_{\mathrm{eff}}}

\title{$N_{\mathrm{eff}}$ Audit (Critical)\\\large Step 2 Audit of Vacuum-Polarization Counting}
\author{Ing. David Jaro\v{s}}
\date{2026-05-10}

\begin{document}
\maketitle

\begin{abstract}
This audit resolves the convention conflict in
\texttt{canonical/n\_eff/step2\_vacuum\_polarization.tex}.
Verdict is explicitly split into two routes:
SU(2)-twist route (A) gives
$N_{\mathrm{helicity}}=2$, $N_{\mathrm{charge}}=2$,
$N_{\mathrm{eff}}=12$, $B_0=8\pi$ at [L1];
the independent loop-counting route (B) gives
$N_{\mathrm{eff}}^{\mathrm{loop}}=3$, $B_0^{\mathrm{loop}}=2\pi$ at [L1],
while the identification with Route A remains [OPEN/MC].
\end{abstract}

\tableofcontents

\section{Two distinct $N_{\mathrm{eff}}$ quantities}

\subsection{Route A: SU(2) twist [L1]}
\[
N_{\rm eff}^{\rm twist} = N_{\rm phases}\times N_{\rm charge}
\times N_{\rm helicity} = 3\times 2\times 2 = 12 \quad [\mathrm{L1}].
\]
Source: \texttt{research\_tracks/quantum\_ubt/su2\_twist\_neff12.tex}.

\subsection{Route B: loop-counting [MC/OPEN]}
\[
N_{\rm eff}^{\rm loop}
\]
is defined by the direct one-loop $\beta$-function count of $S[\Theta]$.
For the minimal scalar action audited here, the charged complex-scalar loop gives
\[
N_{\rm eff}^{\rm loop} = 3 \quad [\mathrm{L1}].
\]
The identification
\[
N_{\rm eff}^{\rm loop} = N_{\rm eff}^{\rm twist}
\]
is \textbf{not proved}. Status: \textbf{[MC/OPEN, frozen]}.

\begin{statusbox}
Route A (SU(2) twist): \textbf{[L1]}\\
Route B (loop): \textbf{[L1 for $N_{\rm eff}^{\rm loop}=3$]}\\
Identification A=B: \textbf{[NO PROOF --- open problem]}
\end{statusbox}

\section{Scope and hard rule}
The audited target is the one-loop coefficient entering
\begin{equation}
\frac{1}{\alpha(\mu)}=\frac{1}{\alpha(\mu_0)}+\frac{B}{2\pi}\ln\frac{\mu}{\mu_0}.
\end{equation}
Every multiplicity factor must be tied either to an explicit loop trace/diagram
or to a non-redundant algebraic counting statement.

\section{Step A --- Nature of the $\Theta$ field in the audited loop}
The minimal canonical action used in the audited chain is
\begin{equation}
S_{\mathrm{kin}}[\Theta]=\int d^4x\,d\psi\;\mathrm{Sc}\!\left[(D_\mu\Theta)^\dagger(D^\mu\Theta)\right],
\end{equation}
and its reduced mode action is of scalar-QED form
$\sum_k |D_\mu\Phi_k|^2$.
No explicit fermionic kinetic term $\bar\Psi i\gamma^\mu D_\mu\Psi$ appears in the
audited derivation path.

\begin{remark}
Therefore the directly audited loop normalization is the complex-scalar loop
normalization, not a Dirac-fermion trace normalization.
\end{remark}

\section{Step B --- Origin of $N_{\mathrm{helicity}}$}
In standard perturbation theory, Passarino--Veltman reduction reduces tensor
integrals to scalar integrals; it does not itself generate a fermionic spinor
multiplicity. Spin multiplicity in a fermion loop comes from explicit Dirac
traces.

For this audited scalar action, the loop coefficient must be taken from the
complex-scalar one-loop vacuum polarization baseline (Peskin \& Schroeder scalar
QED normalization), not from a Dirac trace argument unless a fermionic action is
explicitly introduced and derived.

\textbf{Audit verdict for Step B} (updated per \texttt{canonical/n\_eff/nhelicity\_derivation.tex}):
\begin{itemize}
  \item Three candidate derivation paths for $N_{\mathrm{helicity}}=2$ have been
  examined explicitly:
  \begin{itemize}
    \item \textbf{Path A} (Dirac trace): inapplicable to $S_{\mathrm{kin}}[\Theta]$
    which is a scalar action with no Dirac fermion term.
    \item \textbf{Path B} ($\mathrm{Mat}(2,\mathbb{C})$ off-diagonal components):
    produces $N_{\mathrm{charge}}=2$, not $N_{\mathrm{helicity}}=2$;
    adding it separately is double-counting.
    \item \textbf{Path C} (Weyl decomposition from $\mathbb{C}\otimes\mathbb{H}$):
    physically motivated but requires a derived fermionic effective action
    not present in $S_{\mathrm{kin}}$.
  \end{itemize}
  \item The loop-safe value from $S_{\mathrm{kin}}[\Theta]$ is $N_{\mathrm{helicity}}=1$.
  \item $N_{\mathrm{helicity}}=2$ is \textbf{not proved} as an independent multiplier.
  \item Status: \textbf{OPEN / [MC]} — requires derivation of fermionic effective action.
  \item Dedicated follow-up derivation file:
  \texttt{research\_tracks/quantum\_ubt/fermionic\_action\_derivation.tex}.
\end{itemize}

\section{Step C --- Origin of $N_{\mathrm{charge}}$}
For one \emph{complex} scalar field, particle/antiparticle content is already
encoded in the standard complex-field loop normalization. Adding a separate
$\times 2$ charge factor on top of that same normalization is double counting
unless a distinct non-overlapping degree of freedom is demonstrated.

\textbf{Audit verdict for Step C:}
\begin{itemize}
  \item An extra explicit $N_{\mathrm{charge}}=2$ is \textbf{not justified} in the
  currently audited one-loop normalization.
  \item Status: \textbf{not accepted as an independent multiplier} in Step 2.
\end{itemize}

\section{Step D --- QED-limit normalization}
Using the normalization employed across the current alpha track,
\begin{equation}
B_0=\frac{2\pi}{3}\,N_{\mathrm{eff}}.
\end{equation}
With one complex scalar baseline this gives $B_0^{(1)}=2\pi/3$.

Under the audited minimal-scalar interpretation:
\begin{equation}
N_{\mathrm{eff}}^{\mathrm{loop}}=N_{\mathrm{phases}}=3
\quad\Rightarrow\quad
B_0^{\mathrm{audit}}=2\pi.
\end{equation}
For the stationarity equation used in the alpha track,
\begin{equation}
B=\frac{2n_*}{\ln n_*+1},
\end{equation}
this yields
\begin{equation}
n_*(B_0=2\pi)\approx 10.54,
\end{equation}
which is not 137.

\section{Route A (SU(2)-twist): closed at [L1]}
The Scherk--Schwarz SU(2)-twisted charged sector is tracked as a distinct route
from the scalar-loop baseline. In this route:
\[
N_{\mathrm{phases}}=3,\quad
N_{\mathrm{charge}}=2\ (\text{charged complex fields } b,c),\quad
N_{\mathrm{helicity}}=2.
\]
Therefore
\[
N_{\mathrm{eff}}=N_{\mathrm{phases}}N_{\mathrm{charge}}N_{\mathrm{helicity}}
=3\times2\times2=12,\qquad
B_0=\frac{2\pi}{3}N_{\mathrm{eff}}=8\pi.
\]
Route-A verdict: \textbf{[L1]}.

\section{Route B (loop-counting, independent): frozen [OPEN/MC]}
The independent loop-counting closure is kept separate from Route A.
In this branch, the directly audited scalar loop gives
\[
N_{\mathrm{eff}}^{\mathrm{loop}} = 3
\]
at [L1].  The open blocker is the identification of this direct loop count
with the SU(2)-twist multiplicity $N_{\mathrm{eff}}^{\mathrm{twist}} = 12$.
That identification is explicitly frozen at
\[
\text{status}=\textbf{[OPEN/MC]}.
\]

\section{Consistency check for Task 3 (Mat$(2,\mathbb{C})$, SU(2) left action, chirality)}
The algebraic isomorphism $\mathbb{C}\otimes\mathbb{H}\cong\mathrm{Mat}(2,\mathbb{C})$
and left SU(2) matrix action are internally consistent with the gauge-sector
chirality chain (including the conditional structure in
\texttt{canonical/chirality/step3\_gap\_C1\_resolution.tex}).

\begin{remark}
However, \emph{left matrix action} and \emph{physical chiral fermion loop counting}
are not automatically equivalent in the audited scalar loop. A separate derived
fermionic effective action is required to upgrade $N_{\mathrm{helicity}}$ from
[MC]/open to proved.
\end{remark}

\section{Final verdict table}
\begin{center}
\begin{tabular}{llll}
\toprule
Factor & Value (audit) & Justification & Status \\
\midrule
$N_{\mathrm{phases}}$ & 3 & $\dim_{\mathbb{R}}\mathrm{Im}(\mathbb{H})=3$ & [L0] $\checkmark$ \\
$N_{\mathrm{helicity}}$ & 2 (SU(2) twist) & Anti-periodic off-diagonal KK sector in Scherk--Schwarz route & [L1] \\
$N_{\mathrm{charge}}$ & 2 (claimed) & Extra factor double-counts complex-field baseline unless redefined & OPEN / rejected as independent factor \\
$N_{\mathrm{eff}}^{\mathrm{twist}}$ & 12 & $3\times2\times2$ via SU(2)-twisted charged sector & [L1] \\
$N_{\mathrm{eff}}^{\mathrm{loop}}$ & 3 & Direct scalar-loop count from $S_{\mathrm{kin}}[\Theta]$ & [L1] \\
$B_0^{\mathrm{twist}}$ & $8\pi$ & $B_0=(2\pi/3)N_{\mathrm{eff}}^{\mathrm{twist}}$ & [L1] \\
$B_0^{\mathrm{loop}}$ & $2\pi$ & $B_0=(2\pi/3)N_{\mathrm{eff}}^{\mathrm{loop}}$ & [L1] \\
$N_{\mathrm{eff}}^{\mathrm{loop}}=N_{\mathrm{eff}}^{\mathrm{twist}}$ & not proved & Identification between distinct routes & [MC/OPEN frozen] \\
\bottomrule
\end{tabular}
\end{center}

\section{Impact on alpha track}
This audit does \textbf{not} derive $B\approx 46.3$ from $S[\Theta]$.
It weakens the prior $N_{\mathrm{eff}}=12\Rightarrow B_0=8\pi$ claim by moving the
multipliers $N_{\mathrm{helicity}}$ and explicit $N_{\mathrm{charge}}$ out of proved
status in the currently audited loop derivation.


\begin{statusbox}
$N_{\mathrm{helicity}} = 2$ z SU(2) twist: \textbf{[L1]}\\
$N_{\mathrm{eff}}^{\rm twist} = 12$ a $B_0^{\rm twist}=8\pi$ v SU(2)-twisted větvi: \textbf{[L1]}\\
$N_{\mathrm{eff}}^{\rm loop} = 3$ a $B_0^{\rm loop}=2\pi$ v přímém scalar-loop auditu: \textbf{[L1]}\\
Identifikace twist = loop: \textbf{[OPEN/MC, frozen with blocker note]}\\
Alpha je odvozena: \textbf{NE}
\end{statusbox}

\begin{remark}[Which $N_{\rm eff}$ enters the integer-137 claim]
\label{rem:neff_alpha_dependency}
The conditional result $n^*(B_{\rm phenom}) = 137$ uses
$N_{\rm eff}^{\rm twist} = 12$ [L1] via the A\_PRIME route
(\texttt{canonical/alpha/ALPHA\_MASTER\_STATUS.md}).
The independent loop count $N_{\rm eff}^{\rm loop} = 3$ [L1] gives
$B_0^{\rm loop} = 2\pi$, yielding $n^* \approx 10.5 \neq 137$.
Therefore the T3\_ALPHA companion note is conditional on the
SU(2)-twist route specifically. The identification
$N_{\rm eff}^{\rm twist} = N_{\rm eff}^{\rm loop}$ is not required
for the integer-137 conditional result; it would only be required for
a first-principles derivation of $B_0$ from the scalar effective action.
Status of identification: \textbf{[OPEN/MC frozen]}.
\end{remark}

\end{document}
